% Propietario conservado: /Users/ruben/Documents/New project/output/REVISION_ARGUMENTAL_APERTURAS_CIERRES_20260915/VII/delivery/MANUSCRIPT_VII_ES_EN_COMPLETE/source_es/manuscrito/30c_composicion_corriente_conexion.tex
% Artículo VII, incorporación aditiva, revisión 05.
% Arquitectura autoral: energía de memoria APP--TRIT--TPK, S01.
% Antecedente recuperado: diferencial y refinamiento de 30b.
% Formalización reunida: nota COMPOSICION_VARIACIONAL_TRIT_CONEXION.md.
% Se conserva íntegra 30b; su representante antihermítico es la
% restricción unitaria del diferencial que aquí se desarrolla.
% Esta sección evalúa la corriente de una acción y una realización
% declaradas. No selecciona retrospectivamente su campo ni su peso.

\section{Realization of the current in connection coordinates}
\label{vii:vii:sec:composicion-corriente-conexion}

Memory energy assigns each route a response preserving its
incidences and the order of transports. Inserting it into a geometric
action requires its components with respect to connection
variations. This step is obtained by differentiating the transport that already realizes
the route: the connection current is the composition of that differential
with the energy covector of the preceding section.

The elliptic, parabolic and hyperbolic regimes of TRIT require
preserving the complete differential before restricting it to a
particular representation. Unitary rotations constitute a
domain of that construction. The following extension also preserves
the dilation and transport directions admitted by the geometric
realization used.

\subsection{Differential on invertible transports}

The fibers, fields and weight of
\eqref{vii:vii:eq:energia-memoria-discreta} are retained. We now consider invertible
transports and differentiable families with
\(\dot U_a=A_aU_a\), where \(A_a\) belongs to the tangent space
of the chosen representation. We write
\[
 v_a=U_af_{s(a)},\qquad r_a=f_{t(a)}-v_a,\qquad
 q=\mathsf W_mr,\qquad E_m=\langle r,\mathsf W_mr\rangle.
\]

\begin{proposition}[Complete representative of the differential]
\label{vii:vii:prop:diferencial-invertible}
At fixed field and weight, the following holds
\begin{equation}
 \mathrm d_UE_m(A)=-2\operatorname{Re}\sum_a
       \langle q_a,A_av_a\rangle
 =\sum_a\operatorname{Re}\operatorname{Tr}(\mathcal G_a^*A_a),
 \qquad \mathcal G_a=-2q_av_a^*.
 \label{vii:vii:eq:gradiente-transporte-completo}
\end{equation}
The anti-Hermitian and Hermitian projections of this representative are
\begin{equation}
 \frac{\mathcal G_a-\mathcal G_a^*}{2}
   =v_aq_a^*-q_av_a^*=\mathcal J_a,
 \qquad
 \frac{\mathcal G_a+\mathcal G_a^*}{2}
   =-(q_av_a^*+v_aq_a^*).
 \label{vii:vii:eq:proyecciones-gradiente-memoria}
\end{equation}
Restriction to anti-Hermitian generators recovers exactly
\eqref{vii:vii:eq:representante-corriente-discreta}.
\end{proposition}

\begin{proof}
The identity \(\dot r_a=-A_av_a\) and self-adjointness of the weight give
\(\dot E_m=2\operatorname{Re}\langle q,\dot r\rangle\).
Moreover,
\[
 \operatorname{Tr}(\mathcal G_a^*A_a)
 =-2\operatorname{Tr}(v_aq_a^*A_a)=-2q_a^*A_av_a.
\]
The two projections are obtained by taking the adjoint. For the real
Hilbert--Schmidt inner product, the Hermitian and anti-Hermitian subspaces are
orthogonal; therefore the Hermitian part has zero contraction with
an anti-Hermitian generator.
\end{proof}

The formula allows weights coupling edges: their action is calculated in
\(q=\mathsf W_mr\) before separating the contributions of each
incidence. If field or weight vary, the complete expression
\eqref{vii:vii:eq:variacion-completa-memoria} is retained, also valid for the
invertible transports considered here.

\subsection{Quadratic representations of TRIT}

In a two-dimensional matrix realization of
\(\mathbb A_\tau\), one may take
\[
 J_{+}=\begin{pmatrix}0&1\\-1&0\end{pmatrix},\qquad
 J_{-}=\begin{pmatrix}0&1\\1&0\end{pmatrix},\qquad
 J_{0}=\begin{pmatrix}0&1\\0&0\end{pmatrix}.
\]
Direct multiplication gives
\[
 J_+^2=-I,\qquad J_-^2=I,\qquad J_0^2=0.
\]
Their exponentials realize, respectively, the elliptic,
hyperbolic and parabolic regimes. The first is orthogonal for the fixed
positive inner product. The second preserves the form
\(B=\operatorname{diag}(-1,1)\), since
\(J_-^{\mathsf T}B+BJ_-=0\). The third represents the algebra
of dual numbers. These representations fix examples of the quadratic
types; their use on a screen of another dimension preserves the
corresponding representation map.

A calculation distinguishes the components of the differential. On an edge,
take \(U=I\), \(f_s=(1,1)^{\mathsf T}\),
\(f_t=2f_s\) and \(\mathsf W=I\). Then \(q=v=f_s\),
\(\mathcal J=0\), whereas
\begin{equation}
 \mathrm dE(J_-)=-4,\qquad \mathrm dE(J_0)=-2.
 \label{vii:vii:eq:control-regimenes-corriente}
\end{equation}
The complete representative preserves both responses. The distinction
identifies the space of variations before performing the material
contraction.

\subsection{Composition of routes and oriented inversion}

\begin{proposition}[Contragredient transport of the differential]
\label{vii:vii:prop:pullback-invertible}
For \(U_\gamma=U_N\cdots U_1\), let
\(B_k=U_N\cdots U_{k+1}\) and \(B_N=I\). We have
\[
 \dot U_\gamma U_\gamma^{-1}
   =\sum_kB_kA_kB_k^{-1}.
\]
The representative \(\mathcal G_\gamma\) in the final fiber induces
at incidence \(k\)
\begin{equation}
 \mathcal G_k=B_k^*\mathcal G_\gamma(B_k^{-1})^*.
 \label{vii:vii:eq:transporte-contragrediente}
\end{equation}
\end{proposition}

\begin{proof}
The product rule, followed by multiplication by
\(U_\gamma^{-1}\), proves the first equality. Cyclicity of the
trace gives
\[
 \operatorname{Re}\operatorname{Tr}
  (\mathcal G_\gamma^*B_kA_kB_k^{-1})
 =\operatorname{Re}\operatorname{Tr}
  ((B_k^*\mathcal G_\gamma(B_k^{-1})^*)^*A_k).
\]
For unitary \(B_k\), the conjugation of
\ref{vii:vii:prop:composicion-respuesta} is recovered.
\end{proof}

\begin{proposition}[Inversion by energy congruence]
\label{vii:vii:prop:inversion-congruencia}
For a local energy per edge, inversion is represented by
\begin{equation}
 U_{\bar a}=U_a^{-1},\qquad
 r_{\bar a}=-U_a^{-1}r_a,\qquad
 W_{\bar a}=U_a^*W_aU_a.
 \label{vii:vii:eq:inversion-peso-general}
\end{equation}
It preserves energy. If \(W_a\) remains fixed,
\[
 A_{\bar a}=-U_a^{-1}A_aU_a,\qquad
 \dot W_{\bar a}=U_a^*(A_a^*W_a+W_aA_a)U_a.
\]
The complete variation of the inverted energy coincides with the original.
\end{proposition}

\begin{proof}
Substitution into \(r_{\bar a}^*W_{\bar a}r_{\bar a}\)
produces \(r_a^*W_ar_a\). The derivative of the inverse is
\(\dot U_a^{-1}=-U_a^{-1}A_a\); that of the weight is obtained by
the product rule. Equality of energies holds for the entire
family, so their derivatives coincide, including the term
\(\langle r_{\bar a},\dot W_{\bar a}r_{\bar a}\rangle\).
\end{proof}

Congruence preserves positivity for every invertible transport.
In the unitary representation it reduces to the inversion rule of
\ref{vii:vii:prop:inversion-respuesta}. The sum uses one orientation
per edge, or half the sum over both orientations.

\subsection{Evaluation of current components}

Let \(\theta=(\theta_b)\) be a real variation of the connection
coordinates of the realization used. The effective differential of
transport defines the operators
\begin{equation}
 A_a(\theta)=\dot U_aU_a^{-1}
     =\sum_bB_{ab}\theta_b.
 \label{vii:vii:eq:diferencial-conexion-transporte}
\end{equation}
For a finite composition, their contributions are obtained through
\eqref{vii:vii:eq:transporte-contragrediente}. The operators
\(B_{ab}\) are the derivatives of the connection representation that
realizes the route.

\begin{theorem}[Connection current of the memory action]
\label{vii:vii:thm:corriente-conexion-evaluada}
For a material term \(S_{\rm mem}=\mathfrak a_m E_m\),
with \(\mathfrak a_m\) its normalization on the action line,
at fixed field, weight and normalization we have
\begin{equation}
 \delta S_{\rm mem}=-\frac12\sum_b\theta_b s_b,\qquad
 s_b=4\mathfrak a_m\operatorname{Re}
        \sum_a\langle q_a,B_{ab}v_a\rangle.
 \label{vii:vii:eq:componentes-corriente-conexion}
\end{equation}
If the cellular pairing between the connection and current bases
is an invertible matrix \(M\), defined by
\(\delta S_{\rm mem}=-\tfrac12\theta^{\mathsf T}M\sigma\),
the current coordinates are
\begin{equation}
 \sigma=M^{-1}s.
 \label{vii:vii:eq:coordenadas-corriente-material}
\end{equation}
\end{theorem}

\begin{proof}
Substitute \eqref{vii:vii:eq:diferencial-conexion-transporte} into
\eqref{vii:vii:eq:gradiente-transporte-completo}, multiply by
\(\mathfrak a_m\), and collect the coefficients of
\(\theta_b\). The variational convention \(-1/2\) fixes the
factor \(4\). Comparing both expressions for every variation
\(\theta\) gives \(M\sigma=s\); invertibility of the pairing
proves existence and uniqueness of these coordinates.
\end{proof}

The formula explicitly realizes the composition mentioned in
\ref{vii:vii:subsec:realizacion-respuesta-discreta}. To evaluate each
component, the field is transported, the residual is calculated, the
weight is applied and the result is contracted with the connection variation of the same realization.
The pairing \(M\) preserves the orientation and volume of cells.
A diagonal pairing gives the corresponding division by their
volumes; other bases preserve the complete matrix.

If the connection also enters \(f\), \(\mathsf W_m\)
or \(\mathfrak a_m\), we apply
\eqref{vii:vii:eq:variacion-completa-memoria} and add
\(E_m\delta\mathfrak a_m\). Contributions from all
material terms are summed before composing the Cartan response.

\begin{corollary}[Compatibility of the current with refinement]
\label{vii:vii:cor:refinamiento-corriente-conexion}
For the compatible families of
\ref{vii:vii:thm:refinamiento-respuesta}, let \(P_m\) be the fixed map
taking connection variations from level \(m\) to level
\(m+1\). If the memory actions coincide on those refined fields
and transports, their covectors satisfy
\[
 s_m=P_m^{\mathsf T}s_{m+1},\qquad
 M_m\sigma_m=P_m^{\mathsf T}M_{m+1}\sigma_{m+1}.
\]
\end{corollary}

\begin{proof}
Differentiating the equality of actions in the direction
\(P_m\theta\) gives
\(-\tfrac12\theta^{\mathsf T}s_m
 =-\tfrac12\theta^{\mathsf T}P_m^{\mathsf T}s_{m+1}\).
The arbitrariness of \(\theta\) and the cellular representations
of both covectors prove the two identities. Composition of
these maps iterates the equality to any finite number of refinements.
\end{proof}

\subsection{Insertion into the Cartan equation}

Section~\ref{vii:vii:sec:corriente-respuesta-cartan} normalizes the
material three-form through
\(\delta_\omega S_m=-\tfrac12\int\delta\omega^{IJ}\wedge
\sigma_{IJ}\). The preceding theorem provides its cellular
components for the declared memory action and pairing.
The Cartan response uses that current together with the co-tetrad,
the Barbero--Immirzi parameter and the gravitational coupling of the
same realization.

The material dependence on contorsion determines how that equation
is solved. For affine matter, the current is independent
of contorsion, and the linear inversion and quadratic elimination
proved in Section~\ref{vii:vii:sec:corriente-respuesta-cartan} apply.
For a holonomy
action, the differential is evaluated at the current connection, and
this dependence is retained when solving the equation. For example, an edge with
\(f_s=f_t=1\), unit weight, and \(U(t)=e^{it}\) has
\[
 E(t)=2-2\cos t,\qquad E'(t)=2\sin t.
\]
Here the current depends on \(t\). The complete formula retains
that dependence instead of replacing it by its value at a particular
connection.

Finally, the gravitational density is obtained from the metric
variation of the effective action, with its field and normalization.
The representation of the current, the contortion it produces and the
resulting density are the successive operations of that composition.
The subsequent dilution law transports the amplitude thus evaluated in
the corresponding material section.
